\newpage
\phantomsection
\label{prf:euclid-i-3}
\LRAProofFor{thm:euclid-i-3}

\begin{remark*}[Return]
\hyperref[thm:euclid-i-3]{Return to Theorem}
\end{remark*}

\begin{theorem*}[Euclid I.3: Cutting off an equal segment]
Given two unequal finite straight line segments, there exists a construction
which cuts off from the greater a straight line segment equal to the less.
\end{theorem*}

\begin{proof}[Professional Standard Proof]
\LRAProofBodyStart
Let $AB$ be the greater of the two given unequal segments and let $CD$ be the
less. By Euclid I.2, construct from the point $A$ a segment $AD$ equal to
$CD$. With center $A$ and distance $AD$, describe a circle meeting $AB$ at
$E$.

Since $AE$ and $AD$ are radii of the same circle, $AE=AD$. Since $AD=CD$ by
construction, $AE=CD$ by Common Notion I. Therefore the segment $AE$ has been
cut off from the greater segment $AB$ equal to the less segment $CD$.
\end{proof}

\begin{proof}[Detailed Learning Proof]
\LRAProofBodyStart
\begin{center}
\begin{tikzpicture}[scale=1.2, line cap=round, line join=round]
  \definecolor{euclidblue}{RGB}{30,110,190}
  \tikzset{
    equaltick/.style={
      postaction={decorate},
      decoration={markings, mark=at position 0.5 with {
        \draw[line width=0.6pt] (0,-2.4pt)--(0,2.4pt);
      }}
    }
  }

  \coordinate (A) at (0,0);
  \coordinate (E) at (1.55,0);
  \coordinate (B) at (3.7,0);
  \coordinate (C) at (0,-0.9);
  \coordinate (D) at (1.55,-0.9);
  \pgfmathsetmacro{\r}{1.55}

  \draw[gray!50, thin] (A) circle (\r);
  \draw[very thick] (A) -- (B);
  \draw[very thick, equaltick] (C) -- (D);
  \draw[very thick, euclidblue, equaltick] (A) -- (E);

  \foreach \p/\pos in {A/above left,E/above,B/above right,C/below left,D/below right}{
    \fill (\p) circle (0.023);
    \node[\pos] at (\p) {$\p$};
  }

  \node[font=\footnotesize] at ($(A)!0.5!(B)+(0,0.27)$) {greater segment};
  \node[font=\footnotesize] at ($(C)!0.5!(D)+(0,-0.28)$) {less segment};
  \node[euclidblue, font=\footnotesize] at ($(A)!0.5!(E)+(0,0.34)$) {$AE=CD$};
  \node[gray!70, font=\footnotesize] at ($(A)+(-0.50,-1.50)$) {center $A$};
\end{tikzpicture}
\end{center}


The problem is not merely to draw a segment equal to the shorter one; Euclid
I.2 already does that. The problem is to place such a segment on the greater
given segment. Starting at the endpoint $A$ of the greater segment $AB$, copy
the shorter segment $CD$ to form $AD$. A circle centered at $A$ with radius
$AD$ records exactly that copied length in every direction from $A$.

Where this circle meets the greater segment $AB$, call the point $E$. Then
$AE$ is a radius of the circle and so is equal to $AD$. Since $AD$ was copied
from $CD$, the cut-off segment $AE$ is equal to $CD$.
\end{proof}

\begin{remark*}[Proof structure]
Euclid I.3 depends on I.2 as a length-copying lemma. Once the shorter segment
has been copied from the endpoint of the greater one, Postulate III turns that
copied length into the cutoff point on the greater segment.
\end{remark*}

\begin{dependencies}
\begin{itemize}
  \item \hyperref[def:euclid-straight-line]{Straight line}
  \item \hyperref[def:euclid-circle]{Circle}
  \item \hyperref[thm:euclid-i-2]{Euclid I.2}
  \item \hyperref[ax:euclid-postulate-describe-circle]{Euclid Postulate III}
  \item \hyperref[ax:euclid-common-notion-things-equal-to-same]{Euclid Common Notion I}
\end{itemize}
\end{dependencies}

\clearpage
